\documentclass[10pt]{beamer}
\usetheme{Madrid}
\usepackage[T1]{fontenc}
\usepackage[utf8]{inputenc}
\usepackage[brazil]{babel}
\usepackage{lmodern}
\usepackage{hyperref}
\usepackage{listings}
\usepackage{xcolor}
\setbeamertemplate{navigation symbols}{}

\lstset{
  language=C++,
  basicstyle=\ttfamily\small,
  keywordstyle=\color{blue},
  commentstyle=\color{gray},
  stringstyle=\color{red},
  breaklines=true,
  showstringspaces=false,
}

\title{Estrutura Arduino}
\subtitle{Curso de Microcontroladores}
\author{}
\date{}

\begin{document}

\begin{frame}
  \titlepage
\end{frame}

\begin{frame}{O que é um sketch?}
  \begin{itemize}
    \item O sketch é o programa que roda no Arduino.
    \item Ele começa em `setup()` e continua em `loop()`.
    \item `setup()` é executado apenas uma vez após ligar a placa.
    \item `loop()` repete muitas vezes por segundo.
  \end{itemize}
\end{frame}

\begin{frame}[fragile]{Estrutura básica}
\begin{lstlisting}
const int LED_PIN = 13;

void setup() {
  pinMode(LED_PIN, OUTPUT);
}

void loop() {
  digitalWrite(LED_PIN, HIGH);
  delay(500);
  digitalWrite(LED_PIN, LOW);
  delay(500);
}
\end{lstlisting}
\end{frame}

\begin{frame}{O que cada linha faz}
  \begin{itemize}
    \item `const int LED_PIN = 13;` define uma constante com o número do pino.
    \item `pinMode()` prepara o pino como saída (`OUTPUT`).
    \item `digitalWrite(pin, HIGH)` liga o pino (5V).
    \item `digitalWrite(pin, LOW)` desliga o pino (0V).
    \item `delay(ms)` pausa o programa em milissegundos.
  \end{itemize}
\end{frame}

\begin{frame}{O pino 13}
  \begin{itemize}
    \item O pino 13 tem um LED incorporado na maioria das placas.
    \item Você não precisa de componentes adicionais para testar.
    \item Mas qualquer pino digital pode ser usado.
    \item Lembre-se de conectar resistor se usar pino externo.
  \end{itemize}
\end{frame}

\begin{frame}{Erros comuns}
  \begin{itemize}
    \item `delay()` bloqueia tudo, use para testes apenas.
    \item Não esqueça o resistor (220-1000 Ohms) com LED externo.
    \item Verifique a placa e porta no Arduino IDE.
    \item Use `const` para valores fixos, fica mais legível.
  \end{itemize}
\end{frame}

\begin{frame}{Próximos passos}
  \begin{itemize}
    \item Teste o código acima na sua placa.
    \item Mude o tempo de `delay()` e observe.
    \item Adicione mais LEDs em pinos diferentes.
    \item Na aula seguinte veremos como evitar `delay()` com máquinas de estado.
  \end{itemize}
\end{frame}

\end{document}
